\documentclass[11pt,a4paper]{article}
\usepackage[utf8]{inputenc}
\usepackage[T1]{fontenc}
\usepackage{amsmath,amssymb,amsthm}
\usepackage{listings}
\usepackage{geometry}
\usepackage{hyperref}
\geometry{margin=2.5cm}

\lstset{
  language=C++,
  basicstyle=\ttfamily\small,
  keywordstyle=\bfseries,
  breaklines=true,
  numbers=left,
  numberstyle=\tiny,
  frame=single,
  tabsize=2
}

\title{IOI 2020 -- Mushrooms}
\author{Solution}
\date{}

\begin{document}
\maketitle

\section{Problem Summary}
There are $n$ mushrooms, each either type A or type B. Mushroom 0 is known to be type A. You can query a function \texttt{use\_machine(x)} that takes a sequence of mushroom indices and returns the number of adjacent pairs in the sequence that are of different types. Determine the total count of type A mushrooms using at most a limited number of queries (around $O(\sqrt{n})$ or $O(n^{2/3})$).

\section{Solution Approach}

\subsection{Basic Observations}
\begin{itemize}
  \item \texttt{use\_machine(\{0, i\})} returns 0 if $i$ is type A, 1 if type B. This identifies one mushroom per query.
  \item To be more efficient, we can batch multiple unknowns into a single query.
\end{itemize}

\subsection{Batch Queries}
If we know $k$ mushrooms of type A (say $a_1, a_2, \ldots, a_k$) and want to classify $k$ unknowns $u_1, \ldots, u_k$, we interleave them:

\texttt{use\_machine(\{$a_1, u_1, a_2, u_2, \ldots, a_k, u_k$\})}

The result tells us the total number of type-changes. Since the known mushrooms are all type A, each transition $a_i \to u_i$ contributes 1 if $u_i$ is B, and each transition $u_i \to a_{i+1}$ contributes 1 if $u_i$ is B. So each unknown $u_i$ (except possibly the last) contributes 0 or 2 to the count, and we can deduce the types.

More precisely, arranging as $[a_1, u_1, a_2, u_2, \ldots]$:
\begin{itemize}
  \item Result $r$. Since A-A transition contributes 0 and A-B or B-A contributes 1:
  \item Each B-type unknown contributes exactly 2 to the count (one transition going in, one going out), except the last unknown which contributes 1.
\end{itemize}

This allows classifying $k-1$ unknowns per query, given $k$ known type-A mushrooms.

\subsection{Strategy}
\begin{enumerate}
  \item Start by identifying the first few mushrooms individually.
  \item Once we have a pool of known type-A (or type-B) mushrooms of size $k$, batch-classify $k$ unknowns per query.
  \item As our pool grows, each query becomes more efficient.
  \item Total queries: approximately $O(\sqrt{n})$.
\end{enumerate}

\section{C++ Solution}

\begin{lstlisting}
#include <bits/stdc++.h>
using namespace std;

int use_machine(vector<int> x); // provided by grader

int count_mushrooms(int n) {
    if (n == 1) return 1;
    if (n == 2) {
        return 1 + (1 - use_machine({0, 1}));
    }

    vector<int> A = {0}; // known type A
    vector<int> B;       // known type B

    // Identify mushrooms 1 and 2
    int r1 = use_machine({0, 1});
    int r2 = use_machine({0, 2});
    if (r1 == 0) A.push_back(1); else B.push_back(1);
    if (r2 == 0) A.push_back(2); else B.push_back(2);

    int countA = A.size();
    int idx = 3;

    while (idx < n) {
        // Use whichever pool is larger
        vector<int>& pool = (A.size() >= B.size()) ? A : B;
        bool using_A = (&pool == &A);

        int k = pool.size();
        int batch = min(k, n - idx); // number of unknowns to classify

        // Build interleaved sequence
        vector<int> seq;
        for (int i = 0; i < batch; i++) {
            seq.push_back(pool[i]);
            seq.push_back(idx + i);
        }
        // Optionally add one more known at the end for the last unknown
        if (batch <= k) {
            // The sequence is [p0, u0, p1, u1, ..., p_{batch-1}, u_{batch-1}]
            // Transitions: p0-u0, u0-p1, p1-u1, ..., p_{batch-1}-u_{batch-1}
            // Total transitions = 2 * (number of B-type unknowns if using_A pool)
            //                      ... well, let's compute carefully.
        }

        int r = use_machine(seq);

        // Decode: seq = [p0, u0, p1, u1, ..., p_{b-1}, u_{b-1}]
        // Transitions between consecutive elements:
        // (p0,u0), (u0,p1), (p1,u1), (u1,p2), ...
        // If using A pool: all p_i are type A.
        // Transition (p_i, u_i): 1 if u_i is B, 0 if A.
        // Transition (u_i, p_{i+1}): 1 if u_i is B, 0 if A.
        // So each u_i contributes 2 if B (or 0 if A), except u_{batch-1}
        // which only has the left transition (p_{batch-1}, u_{batch-1}).
        // Total r = 2 * (# B among u_0..u_{batch-2}) + (u_{batch-1} is B ? 1 : 0)

        // Determine type of last unknown
        int last_contrib = r % 2; // 1 if last unknown is B (when using A pool)
        int rest_B = (r - last_contrib) / 2;

        // But we don't know which specific unknowns (besides last) are B.
        // We only know the total count.

        // Actually, we CAN determine each unknown's type individually if we
        // use a smarter interleaving. Let me reconsider.

        // Alternative: use [p0, u0, p1, u1, ...] and process prefix queries.
        // But that uses too many queries.

        // For the batch approach, we ONLY learn the total count, not individual types.
        // So we add all unknowns to a "pending" list and only count.

        // Correction: we do learn each type! Consider:
        // seq = [p0, u0, p1, u1, p2, u2, ...]
        // r = sum of |type(seq[i]) - type(seq[i+1])| for adjacent pairs
        // The transitions are:
        //   t0 = |A - type(u0)|
        //   t1 = |type(u0) - A|
        //   t2 = |A - type(u1)|
        //   t3 = |type(u1) - A|
        //   ...
        // So t_{2i} + t_{2i+1} = 2 * (u_i is B), for i < batch-1
        // And t_{2*(batch-1)} = (u_{batch-1} is B)
        // But we only get the sum r, not individual t values!

        // So we know: rest_B = total B unknowns among u_0..u_{batch-2}
        // and whether u_{batch-1} is B.

        // We DON'T know which specific ones are B among u_0..u_{batch-2}.

        // For counting purposes, this is sufficient!
        // We just need the total count of A mushrooms, not their identities.

        // Count A types in this batch:
        int batchA, batchB;
        if (using_A) {
            // last unknown: B if last_contrib=1, A if 0
            batchB = rest_B + last_contrib;
            batchA = batch - batchB;
        } else {
            // Using B pool: transitions flip
            // t_{2i} + t_{2i+1} = 2 * (u_i is A), etc.
            int rest_A = (r - last_contrib) / 2;
            int lastA = last_contrib;
            batchA = rest_A + lastA;
            batchB = batch - batchA;
        }

        countA += batchA;

        // We need to grow our pools. We can't add unknowns since we don't know
        // their individual types. So we sacrifice query efficiency for pool growth.

        // Better approach: for the first few, identify individually to grow pools.
        // Then batch for counting.

        // Actually, let's refine: process unknowns ONE AT A TIME when pool is small,
        // and batch when pool is large enough.

        // For simplicity, let me just use individual queries for small pool
        // and batch for large.

        idx += batch;
    }

    // Actually, the above logic mixes batching strategies. Let me rewrite cleanly.
    // (The code below is a clean reimplementation.)

    return countA;
}

// Clean reimplementation:
int count_mushrooms_v2(int n) {
    if (n <= 2) {
        if (n == 1) return 1;
        return 1 + (1 - use_machine({0, 1}));
    }

    vector<int> known[2]; // known[0] = type A indices, known[1] = type B indices
    known[0].push_back(0);

    int countA = 1;
    int idx = 1;

    while (idx < n) {
        // Use the larger pool
        int t = (known[1].size() > known[0].size()) ? 1 : 0;
        int k = known[t].size();

        if (k < 2 || n - idx < 2) {
            // Identify one at a time
            int r = use_machine({known[t][0], idx});
            int type_idx = (r == 0) ? t : (1 - t);
            if (type_idx == 0) countA++;
            known[type_idx].push_back(idx);
            idx++;
        } else {
            // Batch: interleave known[t] with unknowns
            int batch = min(k, n - idx);
            vector<int> seq;
            for (int i = 0; i < batch; i++) {
                seq.push_back(known[t][i % k]);
                seq.push_back(idx + i);
            }
            int r = use_machine(seq);

            // Last unknown: contributes r%2 transitions
            int last_diff = r % 2;
            int rest_diff = (r - last_diff) / 2;
            // rest_diff = number of unknowns (among first batch-1) that differ from type t
            // last_diff = 1 if last unknown differs from type t

            int diff_count = rest_diff + last_diff;
            int same_count = batch - diff_count;

            if (t == 0) {
                countA += same_count;
            } else {
                countA += diff_count;
            }

            // We know the type of the last unknown; add it to pool
            int last_type = last_diff ? (1 - t) : t;
            known[last_type].push_back(idx + batch - 1);

            // For batch-1 unknowns, we know total but not individual types
            // We can't add them to pools. That's fine for counting.

            idx += batch;
        }
    }

    return countA;
}

int main() {
    int n;
    scanf("%d", &n);
    printf("%d\n", count_mushrooms_v2(n));
    return 0;
}
\end{lstlisting}

\section{Complexity Analysis}
\begin{itemize}
  \item \textbf{Query complexity:} Initially, we use $O(1)$ query per mushroom. Once the pool reaches size $k$, each query classifies $k$ mushrooms. The pool grows by 1 per batch, so after $k$ batches, pool has size $k + k$ and we classify $k^2$ mushrooms total. Overall queries $\approx O(\sqrt{n})$.
  \item \textbf{Time complexity:} $O(n)$ for processing all mushrooms.
  \item \textbf{Space complexity:} $O(n)$.
\end{itemize}

\end{document}
